% A09 INDEPENDENT AUDIT ONLY; 2026-10-03.
% Audited P09 v1.0 SHA-256:
% b2366edf1ff71199b7f1d02a58cf828148d9a697dfafc9e07e70a523e4f26680
% All six layer verdicts and OVERALL VERDICT: VERIFIED.
% Required corrections: none. This file does not promote P09's lifecycle.
\documentclass[11pt]{article}
\usepackage[T1]{fontenc}
\usepackage[utf8]{inputenc}
\usepackage{amsmath,amssymb,amsthm}
\usepackage[margin=0.85in]{geometry}
\usepackage{longtable,array,booktabs,enumitem}
\usepackage[hidelinks]{hyperref}
\setlength{\emergencystretch}{3em}
\setlength{\parskip}{4pt}
\setlength{\parindent}{0pt}
\newcommand{\one}{\mathbf{1}}
\newcommand{\R}{\mathbb R}
\newcommand{\C}{\mathbb C}
\newcommand{\U}{\mathcal U}
\newcommand{\Orth}{\R_{\geq0}^n}
\newcommand{\norm}[1]{\lVert#1\rVert}
\newcommand{\lin}[1]{P09 lines #1}
\newcommand{\pass}{\textbf{PASS}}
\newcommand{\verified}{\textbf{VERIFIED}}
\newcommand{\ii}{\mathrm i}
\title{A09: Independent Adversarial Audit\\Transcritical Bifurcation at the Origin}
\author{Directed Cyclic Networks research project}
\date{3 October 2026 \quad Audited version: P09 v1.0}
\begin{document}
\maketitle

\begin{tabular}{@{}p{0.76\textwidth}l@{}}
\toprule
CRITICAL SPECTRUM / REDUCTION VERDICT: & \verified\\
GENERIC TRANSCRITICAL VERDICT: & \verified\\
STABILITY / PHYSICAL-BRANCH VERDICT: & \verified\\
GENERIC CRITICAL-DYNAMICS VERDICT: & \verified\\
DEGENERATE CRITICAL-STABILITY VERDICT: & \verified\\
CODIMENSION / SCOPE VERDICT: & \verified\\
\midrule
OVERALL VERDICT: & \verified\\
\bottomrule
\end{tabular}

\textbf{Required mathematical or wording corrections: none.}
The exact submitted source proves the generic full-system local transcritical
bifurcation and the separate cubic-degenerate result within its stated local
domains. In particular, the full origin at $b=\mu=0$ is locally asymptotically
stable in $\U$, and is not locally exponentially stable. The highest-priority
Lyapunov remainder and quadratic-form arguments are valid as written; their
details are reconstructed below, independently of the developer's checker.

This is an audit verdict, not a lifecycle edit. P09 and every existing project
file remain unchanged. No P10 proof is begun. The original P09 development
chat owns administrative closure and the complete audited project archive.

\section{Exact current-file authentication and audit contract}
\label{sec:auth}
The authoritative input is \path{cyclic_network_project_P09_draft.zip}.
The manifest's CURRENT P09 entry is at lines 693--732, with the exact full-file
hash at line 730. It supersedes historical TODO/not-started records.
The audited file is precisely
\path{cyclic_network_project/proofs/P09_transcritical_bifurcation.tex},
P09 v1.0, dated 3 October 2026: 24,155 bytes and 539 lines.

\begin{description}[leftmargin=0pt,labelwidth=0pt,itemindent=0pt]
\item[Full P09 SHA-256:]
\leavevmode\\\nolinkurl{b2366edf1ff71199b7f1d02a58cf828148d9a697dfafc9e07e70a523e4f26680}
\item[Mathematical body SHA-256, line 17 through EOF:]
\leavevmode\\\nolinkurl{a0830944d46d26de35faf09c5f0e4d9f42d6f10620ca22373f6bcf7031705393}
\item[Original complete draft ZIP SHA-256:]
\leavevmode\\\nolinkurl{61df24eb644771cbb84027447c4a21ba1c1d0e1d44dc8c35cf34a0a0c755c8de}
\end{description}

The two existing draft ZIP copies are byte-identical. The archive passes its
CRC check, contains exactly 143 files under the sole root
\path{cyclic_network_project/}, and every member matches the located
\path{output/p09/cyclic_network_project/} working copy. Thus no earlier
placeholder or reconstructed source is being audited. The current-file check
was completed before accepting any mathematical conclusion.

All ten requested control files were located and authenticated: manifest,
claim registry, frozen theorem package, final novelty verdict, novelty
decisions, dependency map, compendium source, shared macros, canonical model,
and notation/assumptions. Current P01--P05, P07, P08 and their A01--A05, A07,
A08 sources were also located. The additional open-decisions file was
consulted for B03 and NB07. Full hashes, sizes and line counts for these 26
inputs are in \path{A09_support/input_inventory.json}; all 143 archive-member
hashes are in \path{A09_support/archive_member_hashes.json}.

Each current prerequisite body hash matches its header's declared current
body hash. Comparison against the complete P07-audited baseline confirms that
all prerequisite proofs and prior audits are unchanged; the only five changed
members in the P09 draft are P09 itself, manifest, registry, compendium source
and compendium PDF. A01/A02 E01 and A05/A08 E01 are already closed as recorded;
A03/A04/A07 require no correction. A02's restored audit source is present:
its historical source-absence header is explicitly superseded by the current
manifest. A04/A05 historical numerical-provenance qualifications remain
nonblocking and are not silently declared resolved here.

The frozen documents retain historical status prose. Their current scope is
consistent with the CURRENT P09 entry: supporting local branch organization;
no two-parameter unfolding; no new priority claim. Exact-model priority
remains unclosed. This audit performed no new literature survey.

\section{Dependency authentication}
\label{sec:deps}
The frozen mathematical edges remain
\[
 \mathrm{P02},\mathrm{P03},\mathrm{P04}\longrightarrow\mathrm{P09},
 \qquad \mathrm{P01}\text{ inherited}.
\]
The following are the actual imported facts, with their original hypotheses.
\begin{longtable}{@{}p{0.16\textwidth}p{0.77\textwidth}@{}}
\toprule Input & Use and scope checked\\\midrule\endhead
P01 & Smooth local field and continuation on $\U=(-1,\infty)^n$;
nonnegative invariance when restricting full-domain stability to $\Orth$.
The canonical positive tuple satisfies P01's weaker hypotheses.\\
P02 & Exact invariant synchrony and scalar restriction (invariance proposition
and scalar equation); algebraic repeated-root locus (loci proposition).
No global nonsynchronous classification is imported.\\
P03 & Original Jacobian and critical-origin structure; direct differentiation
in this audit verifies the same formula, including additive offset collisions.\\
P04 & Homogeneous circulant coefficients and Fourier spectrum at every
synchronous state $r>-1$, even when negative or nonstationary; uniform
derivative identity; nonlinear implications of a Hurwitz Jacobian or a
positive-real-part eigenvalue.\\
Standard theory & Finite-dimensional real IFT; fixed-parameter and extended
center manifolds; stability reduction with stable complement; strict
Lyapunov/continuation arguments; regular level-set theorem. These are ordinary
mathematical tools, not extra project-module edges.\\
\bottomrule
\end{longtable}
The only external proof-label references in P09 are:
\begin{quote}\raggedright
\path{prop:P01-local}, \path{prop:P02-invariance}, \path{eq:P02-scalar},
\path{eq:P03-origin-jacobian}, \path{eq:P04-synchronous-derivative}, and
\path{prop:P02-loci}.
\end{quote}
Every label exists in the authenticated source.
Inspection of uncited arguments also finds no reliance on P05, P07, P08 or
P10--P18. References to these modules only describe scheduling or exclusions.
In particular, transverse stability uses a finite collection of continuous
eigenvalue formulas, not the P05 selection theorem. \textbf{Dependency verdict:
PASS.}

\section{Layer A: critical spectrum and real transverse reduction}
\label{sec:reduction}
All statements fix $a,k,c,h>0$, $n\ge5$, $1\le q<n$. Set
\[
 A=k+h>0,\quad b=a-A,\quad g_c=a+c,\quad\mu=g_c-g,
 \qquad\frac{d\mu}{dg}=-1.
\]
Restrict $|\mu|<g_c/2$ so that $g>0$. The original field is
\begin{equation}
 F_i=a x_i-kx_i^2+\frac{c x_{i-1}}{1+x_i}
       -\frac{(g_c-\mu)x_i}{1+x_{i+q}}-h x_i x_{i-q}.
 \label{eq:original}
\end{equation}
Its denominators do not vanish on $\U$; it is jointly smooth there.

\subsection{Independent critical differentiation and gap}
At zero, the input numerator derivative contributes $c$ at predecessor
column $i-1$, its denominator derivative vanishes, and the rational loss
contributes $-g$ only at column $i$. Both derivatives of the $h$ product and
the quadratic self-loss derivative vanish. Therefore
\[
 D_xF(0,\mu)=(a-g)I+cP=\mu I+c(P-I),\qquad J_c=c(P-I).
\]
All $q$ terms have vanished before any Fourier argument. A disconnected
step-$q$ graph cannot create extra kernel vectors because the surviving
predecessor shift $P$ is a single $n$-cycle.

For $v^{(m)}_j=n^{-1/2}e^{\ii\theta_m j}$,
$Pv^{(m)}=e^{-\ii\theta_m}v^{(m)}$; hence
\[
 \lambda_m^c=c(e^{-\ii\theta_m}-1),\qquad\theta_m=2\pi m/n.
\]
The $n$ vectors form a unitary basis, so this lists algebraic multiplicities.
Only $m=0$ gives zero, and the characteristic polynomial
$(\lambda+c)^n-c^n$ has derivative $nc^{n-1}\ne0$ there. The equation $Pz=z$
forces every coordinate equal, giving real kernel $\R\one$ and geometric
multiplicity one.

For $m\ne0$, $j=\min(m,n-m)$ satisfies
$1\le j\le\lfloor n/2\rfloor$ and
$2\pi j/n\in[2\pi/n,\pi]$. Monotonicity of cosine yields
\[
 \Re\lambda_m^c=-c(1-\cos\theta_m)\le-\delta_n,
 \qquad\delta_n=c(1-\cos(2\pi/n))>0.
\]
For even $n$, the self-conjugate label $m=n/2$ has eigenvalue $-2c$;
for odd $n$, the upper representative angle is strictly below $\pi$ and
the same inequality holds. Equality at the closest nonzero labels is
allowed. Neither the formula nor its proof contains $q$ or $\gcd(n,q)$.
\lin{54--85}: \pass.

\subsection{A real IFT with explicit local quantifiers}
Let $w^T=n^{-1}\one^T$, $E=\ker w^T$ and $\Pi=I-\one w^T$.
Then $w^T\one=1$, $\Pi^2=\Pi$, $\Pi^T=\Pi$ and
\[
 x\longleftrightarrow(r,y)=(w^Tx,\Pi x),\quad
 (r,y)\longmapsto r\one+y
\]
are inverse real linear maps between $\R^n$ and $\R\times E$.
In particular $\norm{x}_2^2=nr^2+\norm y_2^2$.
Since $w^TP=w^T$, $J_c$ preserves $E$ and $w^TJ_c=0$.

Define $T(r,y,\mu)=\Pi F(r\one+y,\mu)$. It maps to the real space
$E$ and $D_yT(0,0,0)=J_c|_E$. To justify real invertibility without any
hidden complex step, if a real $y\in E$ has $J_cy=0$, the kernel calculation
forces $y=t\one$, and $w^Ty=0$ forces $t=0$. Injectivity on the
finite-dimensional real space $E$ gives surjectivity and a real inverse.
Equivalently, its complexification has no zero eigenvalue; a singular real
restriction would produce a complex kernel vector. Thus P09's
complexification argument is sufficient and correct.

The real IFT gives open neighborhoods $I_r,I_\mu$ of zero and an open
$Y\subset E$ containing zero such that, within $I_r\times Y\times I_\mu$,
\[
 T(r,y,\mu)=0\quad\Longleftrightarrow\quad y=\psi(r,\mu).
\]
These sets can be chosen with $r\one+y\in\U$ throughout. Substitution into
the original field gives
\[
 F(r\one,\mu)=f(r,\mu)\one,\qquad
 f=ar-Ar^2+\frac{(c-g)r}{1+r}
   =\frac{r}{1+r}(\mu+br-Ar^2).
\]
Consequently $T(r,0,\mu)=0$ for \emph{every} $(r,\mu)\in I_r\times I_\mu$.
The fixed neighborhood $Y$ contains zero for every such pair; uniqueness
forces $\psi\equiv0$. No state-dependent choice of $Y$ is involved.
Taking a smaller open state ball whose $(r,y)$ image lies in $I_r\times Y$
and a smaller parameter interval produces the claimed common product in
$(x,\mu)$ coordinates.

Every equilibrium in that product is therefore synchronous, and conversely
it is an equilibrium exactly when
\[
 r(\mu+br-Ar^2)=0.
\]
Multiplication by $1+r>0$ neither removes nor creates a root; $r=-1$ is
excluded by the domain. This proof uses no division by $b$, so covers
$b=0$. It rules out all nearby nonsynchronous equilibria, not merely
smooth nonsynchronous branches. \lin{87--130}: \pass.

\textbf{CRITICAL SPECTRUM / REDUCTION VERDICT: VERIFIED.}

\section{Layer B: generic transcritical geometry and dynamics}
\label{sec:generic}
\subsection{Full-vector derivatives, normalization and orientation}
The scalar restriction gives $f(0,\mu)=0$, $f_r(0,\mu)=\mu$ and
$f_\mu=r/(1+r)$. Thus $f_r(0,0)=0$, $f_{r\mu}(0,0)=1$ and, since
\[
 f(r,0)=br^2-(A+b)r^3+O(r^4),
\]
$f_{rr}(0,0)=2b$. The quadratic Taylor coefficient is
$\tfrac12 f_{rr}=b$, not $2b$.

For $v=\one$ and $w^T=n^{-1}\one^T$, $w^Tv=1$, $J_cv=0$ and
$w^TJ_c=0$. Directly in (\ref{eq:original}),
\[
 (F_\mu)_i=\frac{x_i}{1+x_{i+q}},\quad
 D_xF_\mu(0)=I,\quad F_\mu(0)=0.
\]
Replacing $\mu$ by $g$ reverses the sign, so
$w^TF_{x\mu}v=1$ and $w^TF_{xg}v=-1$.

To check every second directional contribution, substitute $x=tv$ at
$g=g_c$ and differentiate each original summand twice at $t=0$:
\[
\begin{array}{c|ccccc}
\text{summand}&at&-kt^2&ct/(1+t)&-g_ct/(1+t)&-ht^2\\\hline
\text{second derivative}&0&-2k&-2c&2g_c&-2h
\end{array}
\]
Their sum in each of the $n$ components is
$2(g_c-c-k-h)=2b$. Averaging with $w$ cancels the factor $n$ exactly:
$w^TF_{xx}[v,v]=2b$. These are full-vector coefficients, not coefficients
inferred from an unverified scalar reduction. For $b\ne0$, mixed and
quadratic nondegeneracy both hold; at $b=0$ the second one fails.
\lin{134--169}: \pass.

\subsection{Unique local nonzero branch and both slopes}
When $b\ne0$, $Q(r,\mu)=Ar^2-br-\mu$ has $Q_r(0,0)=-b\ne0$.
A second real IFT gives the unique root $r_*(\mu)$ near zero.
$Q(0,\mu)=-\mu$ implies $r_*(\mu)=0$ only at $\mu=0$.
The other algebraic root tends to $b/A\ne0$; for example, choose a state
radius less than $|b|/(3A)$ and shrink the parameter interval so the other
root stays outside. This is a fixed-$b$ statement. No uniform neighborhood
as $b\to0$ is claimed or available.

Implicit differentiation gives
\[
 (2Ar_*-b)r_*'=1,\qquad
 2A(r_*')^2+(2Ar_*-b)r_*''=0.
\]
At zero these yield
\[
 r_*'(0)=-\frac1b,\qquad r_*''(0)=\frac{2A}{b^3},\qquad
 r_*(\mu)=-\frac\mu b+\frac{A\mu^2}{b^3}+O(\mu^3).
\]
Because $d\mu/dg=-1$, $dr_*/dg(g_c)=1/b$. In the $(r,\mu)$ plane
the branch tangents $(0,1)$ and $(-1/b,1)$ are independent. P09 explicitly
restricts its transverse-crossing claim to that plane; it makes no
ambient $(n+1)$-dimensional transversality claim. \lin{174--225}: \pass.

\subsection{Exact center restriction and positive time change}
At fixed $\mu=0$, the invariant synchronous line is tangent to the entire
one-dimensional center space, hence is a smooth local center manifold.
For the extended system with $\dot\mu=0$, $F_\mu(0,0)=0$ implies the
linearization $\operatorname{diag}(J_c,0)$. The synchronous line times the
parameter interval is invariant and tangent to its full two-dimensional
center space. The remaining $n-1$ eigenvalues are strictly stable.

For $b\ne0$, define $u=Ar^2-br=r(Ar-b)$. Its derivative
$2Ar-b$ is nonzero near zero, so $(r,\mu)\mapsto(u,\mu)$ is a local
smooth coordinate change on this center manifold. Since
$\mu+br-Ar^2=\mu-u$,
\[
 \dot u=(2Ar-b)\frac{r}{1+r}(\mu-u)
        =\alpha(r)u(\mu-u),\quad
 \alpha(r)=\frac{2Ar-b}{(Ar-b)(1+r)}.
\]
The identity extends smoothly to $r=0$ because $Ar-b\ne0$ there and
$\alpha(0)=1$. On $|r|<1/2$, $2A|r|<|b|/2$, both $Ar-b$ and
$2Ar-b$ retain the sign of $-b$, while $1+r>0$. Thus $\alpha>0$ and
$d\tau/dt=\alpha(r)$ gives exactly $du/d\tau=u(\mu-u)$.
This is an orbital equivalence preserving forward time on the exact center
restriction. It does not assert a decoupled equation for arbitrary
transverse perturbations. \lin{227--248}: \pass.

With the equilibrium completeness and full stability proof below, the
persistent branch, simple center eigenvalue, nonzero coefficients, distinct
branch tangents, exact center dynamics and stable transverse complement
justify the full-system nondegenerate transcritical classification. This
conclusion does not follow from scalar geometry alone, and P09 does not
use that invalid shortcut.

\textbf{GENERIC TRANSCRITICAL VERDICT: VERIFIED.}

\section{Layer C: full stability exchange and physical branch}
\label{sec:stability}
\subsection{Uniform sign and transverse persistence}
Differentiating the scalar field before imposing equilibrium gives
\[
 f_r=\frac{\mu+br-Ar^2}{(1+r)^2}
       +\frac{r}{1+r}(b-2Ar).
\]
At a nonzero equilibrium $\mu=-r_*(b-Ar_*)$, so
\[
 \lambda_0^{\rm nz}
 =\frac{r_*}{1+r_*}(b-2Ar_*)
 =-\mu\frac{b-2Ar_*}{(1+r_*)(b-Ar_*)}.
\]
For $|r_*|<1/2$ and $2A|r_*|<|b|/2$, both factors involving $b$
have the sign of $b$, and the displayed fraction is positive. Thus
\[
 \operatorname{sgn}\lambda_0^{\rm nz}=-\operatorname{sgn}\mu
\]
exactly, on either side, for either sign of $b$. Its expansion
$-\mu+O(\mu^2)$ corroborates this sign and is not used as the main proof.
\lin{250--270}: \pass.

Independent differentiation of the original field at $r\one$ gives
\begin{align*}
 J(r,\mu)&=dI+CP+GP^{-q}-HP^q,\\
 d&=a-2kr-\frac{cr}{(1+r)^2}-\frac{g_c-\mu}{1+r}-hr,\\
 C&=\frac c{1+r},\qquad G=\frac{(g_c-\mu)r}{(1+r)^2},\qquad H=hr.
\end{align*}
Coincident columns must be added, as in the defining field. Hence
$\lambda_m=d+Ce^{-\ii\theta_m}+Ge^{\ii q\theta_m}-He^{-\ii q\theta_m}$.
These coefficients remain real and smooth at negative $r>-1$; positivity
of $G$ and $H$ is not being assumed. Each nonzero label has limit
$\lambda_m^c$ at $(0,0)$. Intersecting finitely many continuity
neighborhoods gives a common neighborhood where
\[
 \Re\lambda_m(r,\mu)<-\delta_n/2\quad(m\ne0).
\]
This also holds when $b=0$; the proof never divided by $b$.
\lin{272--292}: \pass.

\subsection{Nonlinear conclusions and number of unstable directions}
At the origin, $\lambda_0=\mu$ and
$\Re\lambda_m=\mu-c(1-\cos\theta_m)$; for $|\mu|<\delta_n/2$
all nonuniform modes remain strictly negative. Intersect this interval
with the reduction, branch and gap neighborhoods. For $\mu\ne0$ all
equilibria under discussion are hyperbolic. P04's checked nonlinear
linearization implications apply to the smooth field on its open domain:
\begin{center}
\begin{tabular}{lll}
\toprule
Parameter side & Origin & Local nonzero branch\\\midrule
$\mu<0$, $g>g_c$ & Locally exponentially stable & Unstable\\
$\mu>0$, $g<g_c$ & Unstable & Locally exponentially stable\\
\bottomrule
\end{tabular}
\end{center}
The unstable uniform eigenvalue is real and simple, and all transverse
eigenvalues have negative real parts. Thus the unstable spectral space,
and the local unstable manifold, have real dimension one at each unstable
\emph{off-critical} equilibrium in this table. P09's ``here'' at line 300
has precisely that local proof context. It does not ascribe a positive
eigenvalue to the nonhyperbolic unstable origin discussed later at
$\mu=0$, nor assert that every perturbation of an unstable point escapes.
No revision of this wording is required. \lin{294--303}: \pass.

\subsection{Positive side and actual nonnegative escape}
The relation $\mu=-r_*(b-Ar_*)$, with the second factor retaining the sign
of $b$, proves $r_*>0$ exactly when
$\operatorname{sgn}\mu=-\operatorname{sgn}b$ for nonzero local $\mu$.
Therefore:
\begin{itemize}[leftmargin=*]
\item $b<0$ ($a<k+h$): the positive branch lies on $\mu>0$, or $g<g_c$,
and is locally exponentially stable; the origin is unstable.
\item $b>0$ ($a>k+h$): the positive local branch lies on $\mu<0$, or
$g>g_c$, and is unstable; the origin is locally exponentially stable.
\end{itemize}
This last branch is the small root, not any other positive root separated
from zero for fixed $b>0$.

At each fixed $\mu>0$, $f(r,\mu)/r\to\mu$ as $r\downarrow0$.
Choose an $\varepsilon>0$ inside the local domain with
$f(r,\mu)\ge(\mu/2)r$ on $(0,\varepsilon]$. A solution with arbitrarily
small $r_0>0$ is increasing and must reach $\varepsilon$ by time at most
$2\mu^{-1}\log(\varepsilon/r_0)$: otherwise the integrated inequality
would contradict $r<\varepsilon$. Continuation holds until exit from
the compact scalar interval. The exact lift $r(t)\one$ is nonnegative
throughout. This proves origin instability relative to $\Orth$ by an
admissible trajectory, not merely a positive eigenvector. The fixed escape
radius may depend on $\mu$, as the local instability definition allows.

For $\mu<0$, full-domain exponential stability and P01's nonnegative
invariance restrict to the orthant. A positive unstable branch point is
interior to the orthant; its positive scalar derivative also supplies
small synchronous departures remaining positive until a fixed local exit.
\lin{305--330}: \pass.

\textbf{STABILITY / PHYSICAL-BRANCH VERDICT: VERIFIED.}

\section{Layer D: generic critical-point dynamics and state-space scope}
\label{sec:critical}
At $\mu=0$, $b\ne0$, the invariant-line equation is exactly
\[
 \dot r=\frac{r^2(b-Ar)}{1+r}.
\]
Choose $0<\varepsilon<1/2$ and $A\varepsilon<|b|/2$. For
$0<|r|\le\varepsilon$, $b-Ar$ has the sign of $b$ and magnitude
at least $|b|/2$, whereas $0<1+r\le3/2$. Thus
$|\dot r|\ge |b|r^2/3$. If $b>0$, take $r_0>0$; if $b<0$, take
$r_0<0$. In both cases $s=|r|$ increases and
\[
 \dot s\ge\frac{|b|}{3}s^2,\qquad
 \frac1{s(t)}\le\frac1{s(0)}-\frac{|b|}{3}t
\]
until $s=\varepsilon$. The exit time is at most
$3|b|^{-1}(s(0)^{-1}-\varepsilon^{-1})$. Smooth continuation within
$[-\varepsilon,\varepsilon]\subset(-1,\infty)$ rules out any earlier
failure of existence. These arbitrarily small data exit one fixed
full-state neighborhood because $\norm{r\one}_2=\sqrt n|r|$.

This proves nonlinear instability in the open domain $\U$ in both generic
cases. For $b>0$ the escaping states are nonnegative. For $b<0$ they are
negative, and the proof does not imply relative-orthant instability.
P09 explicitly makes no full relative-orthant critical stability claim in
that latter case; no subsequent sentence enlarges it. Statements of
orthant instability elsewhere have $\mu>0$, or $b>0,\mu=0$, as their
hypotheses. \lin{332--353}: \pass.

\textbf{GENERIC CRITICAL-DYNAMICS VERDICT: VERIFIED.}

\section{Layer E: degenerate geometry, spectra and center stability}
\label{sec:degenerate}
Set $b=0$, hence $a=A=k+h$. The transverse IFT already proved local
completeness without this genericity assumption. The remaining equation is
\[
 f(r,\mu)=\frac{r}{1+r}(\mu-Ar^2).
\]
For $\mu<0$ there is no nonzero real root; for $\mu=0$ only the origin;
for $\mu>0$ the two nonzero roots are exactly $r_\pm=\pm\sqrt{\mu/A}$.
For $0<\mu<A$ both satisfy $r>-1$, and a further parameter shrinkage puts
both in every required reduction and spectral-gap neighborhood. Only the
positive one is nonnegative. There are no hidden nonsynchronous equilibria.

At the critical point,
\[
 f(r,0)=-\frac{Ar^3}{1+r},\qquad
 f_r(0,0)=f_{rr}(0,0)=0,\qquad f_{rrr}(0,0)=-6A\ne0.
\]
``Triple zero'' is the multiplicity of the scalar field's zero, not a
triple eigenvalue. The full Jacobian is still $J_c$ with one simple zero.
At either nonzero root,
\[
 \lambda_0=f_r(r_\pm,\mu)=-\frac{2Ar_\pm^2}{1+r_\pm}<0.
\]
The negative root is strictly greater than $-1$, so its denominator is
positive; the numerator is negative on \emph{both} branches. Combining
with the same transverse gap proves full nonlinear local exponential
stability of both roots in $\U$. The origin is exponentially stable for
$\mu<0$ and unstable for $\mu>0$, with the earlier synchronous positive
escape proving relative-orthant instability. Both nonzero roots occur
on $g<g_c$, because $\mu=g_c-g$. \lin{358--403}: \pass.

\subsection{The center-manifold theorem is applicable as written}
The reduction principle being used is the fixed-equilibrium,
finite-dimensional center-stability principle: for a sufficiently smooth
local system whose linearization splits into a center spectrum on the
imaginary axis and a strictly stable spectrum, with no unstable spectrum,
local asymptotic stability of the equilibrium in the restriction to a
local center manifold implies local asymptotic stability of the full
equilibrium. A smooth invariant manifold through the equilibrium tangent
to the complete center subspace is a valid local center manifold here.

At the fixed tuple $b=\mu=0$, the original field is smooth, the invariant
synchronous line is tangent to $\R\one$, the remaining eigenvalues have
real part at most $-\delta_n$, and no unstable spectrum exists. On that
line,
\[
 \frac{d}{dt}\frac{r^2}{2}=-\frac{Ar^4}{1+r}<0
 \quad(0<|r|<1/2).
\]
The scalar origin is locally asymptotically stable from both sides.
All hypotheses of the stated reduction principle hold. P09 explicitly
fixes $\mu=0$ and excludes the frozen parameter direction from this
stability application. The extended two-dimensional center manifold
used for the generic normal form is not being misused here. The
application in \lin{405--415} is rigorous as written; the direct proof
below is independent confirmation, not a repair of a defective theorem
application.

\section{Highest-priority audit: direct full-system Lyapunov argument}
\label{sec:lyap}
Everything in this section is at the fixed degenerate critical tuple,
with $g=g_c$, $a=A$. Use the real coordinates $x=r\one+y$, $y\in E$.

\subsection{Why the mean remainder is quadratic in the transverse state}
Let $J(r)=D_xF(r\one,0)$. Independently from the original field,
$J(r)=d(r)I+C(r)P+G(r)P^{-q}-H(r)P^q$, even at a nonstationary
synchronous state. Its column sums, as well as row sums, equal
$d+C+G-H=f_r(r,0)$. Therefore
\[
 w^TJ(r)y=f_r(r,0)w^Ty=0\qquad\text{for every }r\text{ and }y\in E.
\]
This is an identity at each $r$, not just at $r=0$. In particular it
eliminates any putative mixed $r\,y$ term in the mean equation.

Choose a compact convex coordinate neighborhood whose image lies in $\U$,
so the entire segment $r\one+t y$, $0\le t\le1$, lies there. Taylor's
formula in the transverse variable gives
\begin{align*}
 F(r\one+y,0)&=F(r\one,0)+J(r)y\\
 &\quad+\int_0^1(1-t)D_x^2F(r\one+t y,0)[y,y]\,dt.
\end{align*}
After averaging, only the first and last terms remain. The second
derivatives are bounded uniformly on this compact set; hence
\[
 \dot r=-\frac{Ar^3}{1+r}+R(r,y),\qquad
 |R(r,y)|\le K\norm y_2^2.
\]
The constant $K$ is chosen on an initial neighborhood and kept fixed while
the later radius shrinks. There is no circular choice of $K$ and $\rho$.
\lin{419--426}: \pass.

For an additional independent algebraic check, put $s=1+r$ and
$\langle z_i\rangle=n^{-1}\sum_i z_i$. Expanding the original rational
terms exactly gives
\begin{align}
 R={}&-k\langle y_i^2\rangle-h\langle y_i y_{i-q}\rangle\notag\\
 &+c\left\langle\frac{r y_i^2}{s^2(s+y_i)}
              -\frac{y_{i-1}y_i}{s(s+y_i)}\right\rangle\notag\\
 &-g_c\left\langle\frac{r y_{i+q}^2}{s^2(s+y_{i+q})}
              -\frac{y_i y_{i+q}}{s(s+y_{i+q})}\right\rangle.
 \label{eq:explicitR}
\end{align}
This follows from the exact identity
\[
 \frac{r+p}{s+z}=\frac rs+\frac ps-\frac{rz}{s^2}
      +\frac{rz^2}{s^2(s+z)}-\frac{pz}{s(s+z)}.
\]
All linear terms average to zero because each cyclic shift of $y$ has
zero mean. For $|r|,\norm y_2\le1/4$, $s\ge3/4$ and
$s+y_i\ge1/2$. The two rational coefficients are bounded by $8/9$ and
$8/3$ respectively. Cauchy--Schwarz for each shifted product yields the
explicit sufficient constant
\[
 |R|\le K_0\norm y_2^2,\qquad
 K_0=\frac1n\left[A+\frac{32}{9}(c+g_c)\right].
\]
This extra identity is audit evidence only; P09's bounded-Taylor proof
already suffices. It works when offsets collide or $n,q$ are noncoprime.

\subsection{Exact transverse factorization}
Since $\Pi F(r\one,0)=0$, the fundamental theorem of calculus gives
\[
 \dot y=\Pi\int_0^1D_xF(r\one+t y,0)y\,dt
       =J_cy+B(r,y)y,
\]
where, as an operator from $E$ to $E$,
\[
 B(r,y)=\left.\Pi\int_0^1[D_xF(r\one+t y,0)-J_c]\,dt\right|_E.
\]
Here $\Pi J_cy=J_cy$. This is exact, including $y=0$, and not an
approximate factorization requiring division by $\norm y$. Smoothness
implies $B(0,0)=0$ and continuity in operator norm, so
$\norm{B(r,y)}_2\le\delta_n/2$ on a sufficiently small neighborhood.
\lin{427--434}: \pass.

\subsection{The quadratic-form estimate does not rely on nonnormal spectra}
For a real vector $y$, the skew-symmetric part contributes zero:
\[
 y^TJ_cy=y^TSy,\qquad
 S=\frac{J_c+J_c^T}{2}=c\left(\frac{P+P^{-1}}2-I\right).
\]
The real symmetric matrix $S$ is unitarily diagonalized by the same
Fourier basis, with eigenvalues $c(\cos\theta_m-1)$. On $E$ its
uniform Fourier coefficient is zero; by Parseval,
\[
 y^TSy=\sum_{m=1}^{n-1}c(\cos\theta_m-1)|\widehat y_m|^2
 \le-\delta_n\sum_{m=1}^{n-1}|\widehat y_m|^2
 =-\delta_n\norm y_2^2.
\]
This includes the real singleton when $n$ is even. The bound is sharp on
the real sine/cosine space of labels $1,n-1$. P09's short phrase
``Fourier diagonalization gives'' refers to the valid unitary
diagonalization already established in P04. It is not an inference
from the real parts of the spectrum of an arbitrary nonnormal matrix.
The explicit symmetric-part calculation confirms that no correction is
needed at \lin{434--436}.

It follows that
\[
 y^T\dot y\le-\delta_n\norm y_2^2+\norm B_2\norm y_2^2
                 \le-\frac{\delta_n}{2}\norm y_2^2.
\]

\subsection{Strict decrease, invariant sublevels and continuation}
Choose $0<\rho<1/2$ so that the coordinate box
$|r|,\norm y_2<\rho$ lies in all preceding neighborhoods and
$K\rho\le\delta_n/4$. Then for
$V=(r^2+\norm y_2^2)/2$,
\begin{align*}
 \dot V&=-\frac{Ar^4}{1+r}+rR+y^T\dot y\\
 &\le-\frac{A}{1+\rho}r^4+K\rho\norm y_2^2
                  -\frac{\delta_n}{2}\norm y_2^2\\
 &\le-\frac{A}{1+\rho}r^4-\frac{\delta_n}{4}\norm y_2^2<0
                       \quad\text{if }(r,y)\ne(0,0).
\end{align*}
For negative $r$, $0<1+r\le1+\rho$, so the denominator comparison has
exactly the indicated direction. The term $rR$ was bounded in absolute
value; no sign assumption on $r$ or $R$ is used.

For example take $0<\ell<\rho^2/4$. On the closed sublevel $V\le\ell$,
$|r|,\norm y_2\le\sqrt{2\ell}<\rho$, and
\[
 |x_i|\le|r|+\norm y_2\le2\sqrt\ell<\rho<1.
\]
This is a compact subset of $\U$ and of the strict-estimate neighborhood.
Since $\dot V<0$ on its nonzero boundary, it is positively invariant.
The smooth vector field is bounded and locally Lipschitz near this compact
set, so the local solution continues for all $t\ge0$ while remaining
there. The coordinate norm is equivalent to the full state norm:
\[
 \frac{\norm x_2^2}{2n}\le V\le\frac{\norm x_2^2}{2}.
\]
Small invariant sublevels give Lyapunov stability. Within any such compact
set, the only point where $\dot V=0$ is zero. The strict Lyapunov theorem,
or LaSalle on that set, gives convergence to zero. This proves local
asymptotic stability in the \emph{full} open domain, without presuming
global invariance of $\U$. \lin{437--449}: \pass.

\textbf{Explicit substantive conclusion: the full origin at $b=0$, $\mu=0$
is locally asymptotically stable in $\U$. VERIFIED.}

\subsection{Algebraic decay rules out exponential stability}
On synchrony at the same tuple, a small positive solution decreases and
a small negative solution increases. Uniqueness prevents either from
crossing or reaching the equilibrium in finite time. Monotonicity and
the absence of any other scalar equilibrium force convergence to zero.
For every nonzero such solution,
\[
 \frac{d}{dt}(r^{-2})=-2r^{-3}\dot r=\frac{2A}{1+r(t)}\longrightarrow2A.
\]
Integrating and taking a Cesaro mean gives
$r(t)^{-2}/t\to2A$, hence $|r(t)|\sim(2At)^{-1/2}$. Sign preservation
makes this valid from either side. If the full origin were locally
exponentially stable, these exact full-system solutions would satisfy
$\sqrt n|r(t)|\le M e^{-\gamma t}\sqrt n|r_0|$ for some
$M,\gamma>0$. Their algebraic decay contradicts that estimate.
No rate for every nonsynchronous trajectory is inferred.
\lin{451--459}: \pass.

\textbf{DEGENERATE CRITICAL-STABILITY VERDICT: VERIFIED.}

\section{Layer F: codimension, terminology and scope}
\label{sec:scope}
In the open positive parameter space with coordinates $(a,k,c,g,h)$,
set $H_1=a-k-h$ and $H_2=g-a-c$. Their gradients are
\[
 \nabla H_1=(1,-1,0,0,-1),\qquad
 \nabla H_2=(-1,0,-1,1,0).
\]
A linear dependence has second coefficient zero by the $g$ coordinate,
then first coefficient zero by the $k$ coordinate. Equivalently the
minor in columns $k,g$ is $-1$. Rank two holds everywhere. The common
zero set is nonempty, since any $k,c,h>0$ give $a=k+h$, $g=a+c$.
The regular level-set theorem gives a smooth codimension-two
submanifold for each fixed admissible $(n,q)$.

This is rank in physical parameter space, not algebraic multiplicity of
an eigenvalue. At every point on the locus, $J_c=c(P-I)$ still has exactly
one simple zero. P09 expressly excludes a double-zero spectral degeneracy;
there is no Bogdanov--Takens assertion.

Verified P02 gives
\[
 g_f=a+c+\frac{(a-A)^2}{4A},\qquad r_f=\frac{a-A}{2A}.
\]
As $A>0$, $g_f=g_c$ iff $a=A$, at which $r_f=0$. The repeated root is
positive only when $a>A$. Thus the meeting point is the origin and the
statement is only an algebraic meeting of loci. In fact the two equations
$g=g_c$ and $g=g_f$ themselves meet tangentially: the rank-two proof uses
$b=0$ and $g-g_c=0$, not independent gradients of the two locus graphs.
P09 does not confuse these two rank statements. \lin{462--492}: \pass.

P09's degenerate terminology is consistently ``cubic-degenerate origin'',
``square-root nonzero branches'', and failure of ordinary transcritical
quadratic nondegeneracy. Its occurrences of ``pitchfork'' explicitly
withhold that classification. No $\mathbb Z_2$ symmetry is asserted.
The rational scalar field is not presumed odd. No named codimension-two
unfolding is claimed.

The negative degenerate stable branch is a signed equilibrium in $\U$;
it supplies no second positive stable branch for a physical bistability
theorem. The generic $b>0$ case only analyzes the branch near zero.
P09 neither proves nor claims a positive-fold saddle-node, true
bistability, basin separation, basin geometry, cusp unfolding or
two-parameter unfolding. P10 retains those subsequent tasks within the
frozen scope; B03 still requires its precise local coexistence wording.
Likewise there is no P05 selection theorem, five-node specialization,
DM--MB comparison, Hopf assertion, imaginary-pair transversality,
first Lyapunov coefficient, rotating wave or large-$n$ theorem here.
The real branch-crossing word ``transverse'' is not an imaginary-pair
transversality claim. Neighborhoods may depend on $n$; the gap shrinks
with growing $n$ and no uniform-in-$n$ assertion is made.
\lin{494--539}: \pass.

\textbf{CODIMENSION / SCOPE VERDICT: VERIFIED.}

\section{Independent corroboration, provenance and counterexample search}
\label{sec:checks}
\subsection{Independent checks}
The audit's \path{A09_support/independent_checks.py} was initially authored
from the canonical field before reading the developer checker. It imports
no developer code or output. It uses exact rational series division and
rational identities separately from ordinary NumPy floating-point tests.
Its final recorded run used Python 3.12.14 and NumPy 2.3.5, with fixed random
seed 90261003. The source and all raw results are packaged.

\begin{center}
\begin{tabular}{lr}
\toprule
Independent check class & Count\\\midrule
Exact rational identities and derivative checks & 310\\
Original-field equilibrium / direct-matrix cases & 2,370\\
Generic local slope checks, both orientations & 316\\
Original-field transverse Newton searches & 474\\
Full-field Lyapunov / remainder / transverse-factor samples & 1,580\\
\bottomrule
\end{tabular}
\end{center}
Dimensions are $5,6,7,8,9,10,12,15,16$, with every admissible $q$:
79 network pairs, including 33 noncoprime pairs. This includes $q=1$,
$q=n-1$ and $q=n/2$ whenever allowed. Six positive fixed tuples include
both generic signs, two exactly degenerate tuples, a small coupling $c$,
and $b$ close to zero on both sides. Both signs of $\mu$, critical points,
the generic local branch and both degenerate roots are tested.

The original field is evaluated directly. Jacobian entries are assembled
by adding each summand's derivative to its actual receiving column,
including collisions, and checked against complex-step differentiation.
All direct matrix eigenvalues are paired to the analytic list with
multiplicity retained. Maximum equilibrium residual is
$1.214\times10^{-20}$; maximum scaled spectrum discrepancy is
$4.797\times10^{-15}$; maximum scaled Jacobian discrepancy is
$6.939\times10^{-16}$; maximum relative slope error is
$1.251\times10^{-7}$. These maxima are run diagnostics, not rigorous
enclosures.

The transverse searches fix $(r,\mu)$ and solve the original projected
field in a real orthonormal basis of $E$ from small nonzero initial
transverse states. All converge to numerical zero; the largest final
transverse norm is $1.636\times10^{-19}$. This is only a finite search,
not a proof of IFT uniqueness. The Lyapunov tests independently evaluate
$R$, $\dot V$, the symmetric-part eigenvalues, and a quadrature approximation
of $B$. All satisfy the tested inequalities. Tiny-$b$ eigenvalues may
approach floating-point resolution; their sign is checked by the exact
scalar formula, not falsely certified by a rounded matrix eigenvalue.

\subsection{Developer validation is kept separate}
The developer companion was located and its original checker authenticated:
\begin{center}\small
\texttt{b42af82b97c187b3a61238b7cc2aba9d35e44d210a38ec88421b10d5be2fcf93}
\end{center}
It was copied into the audit support directory and rerun there, leaving
the original untouched. All four result files (exact, numerical, slopes,
summary) reproduced byte-for-byte: 1,260 exact checks, 2,700 matrix cases
and 12 slope checks. The originals, rerun files and hash comparison are
under \path{A09_support/prior_validation/}. Their role is validation
provenance only. This reproduction neither proves P09 nor resolves the
older A04/A05 provenance qualifications.

\subsection{Twelve targeted attempts to break the proof}
\begin{longtable}{@{}p{0.06\textwidth}p{0.26\textwidth}p{0.61\textwidth}@{}}
\toprule No. & Target & Adversarial resolution\\\midrule\endhead
1 & $\psi\equiv0$ & Tested parameter-dependent neighborhoods and uniqueness.
A fixed product contains $y=0$ for every nearby $(r,\mu)$; uniqueness forces
zero. This is analytic, not a numerical observation.\\
2 & Equilibrium completeness & A merely scalar restriction would be inadequate.
The invertible real transverse derivative excludes every nearby
nonsynchronous solution, including nonsmooth hypothetical branches.\\
3 & Branch slope sign & Differentiate $Q$ twice and then use $d\mu/dg=-1$;
both signs of $b$ checked. No missing orientation reversal.\\
4 & Stability exchange & Recomputed all modes, including even singleton and
offset collisions. Only the uniform eigenvalue can change sign locally.\\
5 & Positive branch side & The exact relation $\mu=-r_*(b-Ar_*)$ fixes the
side. An additional root away from zero cannot be substituted for $r_*$.\\
6 & Time factor & Both potential zeros $r=b/A$ and $r=b/(2A)$ are excluded
by a fixed-$b$ neighborhood; $r=-1$ is outside the domain. $\alpha(0)=1$.\\
7 & Generic critical escape & For $b<0$, negative $r$ gives increasing
$|r|$, rather than positive escape. The source preserves the orthant scope.\\
8 & Negative degenerate root & $1+r_->0$ locally and the uniform numerator
is $-2Ar_-^2<0$. No positive-state coefficient signs were imported.\\
9 & Full degenerate stability & Both the fixed-parameter reduction principle
and an independent strict Lyapunov function apply; parameter neutrality
cannot be used to refute the fixed-tuple assertion.\\
10 & Mean remainder & A potential $|r|\norm y$ term vanishes identically:
$w^TJ(r)|_E=0$ at every $r$. Exact formula (\ref{eq:explicitR}) checks this.\\
11 & Quadratic-form bound & Real parts alone would fail for a general
nonnormal matrix. Here the symmetric part has exactly the Fourier cosine
eigenvalues, proving the asserted bound.\\
12 & Codimension & Independent gradients of $b$ and $g-g_c$ have rank two;
the tangential meeting of the repeated-root and intersection graphs is
not mistaken for a transverse intersection.\\
\bottomrule
\end{longtable}
No analytic counterexample survives these checks. No finite computation
is claimed to exclude counterexamples universally; the analytic arguments
in Sections~\ref{sec:reduction}--\ref{sec:scope} supply that justification
under P09's stated local hypotheses.

\section{Complete source-line review}
\label{sec:lines}
Line numbers refer to the authenticated 539-line P09 v1.0 source, not its
PDF or an earlier placeholder. The ledger below covers every line once
in consecutive ranges, including metadata, separators and the closing
checklist. Every substantive range is assessed. All ranges are PASS;
there is no MINOR ISSUE, MAJOR ISSUE, FALSE or UNRESOLVED finding.
The machine-readable companion is \path{A09_support/line_review.json}.

% A09-LINE-LEDGER-START
\begin{longtable}{@{}p{0.11\textwidth}p{0.10\textwidth}p{0.70\textwidth}@{}}
\toprule Lines & Verdict & Finding\\\midrule\endhead
1--15 & \pass & Metadata correctly declares DRAFT, fixed parameter path, direct P02--P04 dependencies, inherited P01 and scope exclusions.\\
16--20 & \pass & Title and status agree with the exact CURRENT manifest; no verification is self-awarded.\\
21--27 & \pass & Canonical tuple and abbreviations are consistent; $b$ is unambiguously local and $A>0$.\\
28--33 & \pass & All five original DM terms are transcribed with correct indices and $g=g_c-\mu$.\\
34--39 & \pass & Joint smoothness, positive $g$ neighborhood, open domain and parameter-dependent locality are correct; no global invariance or uniform-in-$n$ claim.\\
40--49 & \pass & P02 import is the exact invariant restriction with denominator $1+r$; no global completeness is inferred here.\\
50--53 & \pass & Subsection and anchor only; no added claim.\\
54--70 & \pass & Critical spectrum lemma: algebraically simple zero, real kernel, explicit positive gap and arbitrary admissible $q$ are correct.\\
71--80 & \pass & Original differentiation, Fourier phases, characteristic-polynomial derivative and real kernel are independently correct.\\
81--85 & \pass & Cosine comparison covers odd/even dimensions and all paired/singleton nonzero labels.\\
86--97 & \pass & Local equilibrium completeness is stated in state and parameter neighborhoods, including $b=0$.\\
98--107 & \pass & Mean coordinates and orthogonal mean-zero projection give a real linear isomorphism; normalization agrees with scalar amplitude.\\
108--115 & \pass & Projected derivative preserves $E$; complexification has no kernel, which does imply invertibility of the real restriction.\\
116--124 & \pass & Real IFT product and uniqueness imply $\psi\equiv0$; zero belongs to one common transverse neighborhood.\\
125--130 & \pass & Both equilibrium implications are proved; multiplication by positive $1+r$ is harmless; shrinking gives a state/parameter product.\\
131--135 & \pass & Subsection/proposition opening has no unsupported claim.\\
136--141 & \pass & Scalar derivative values and mixed derivative are correct.\\
142--154 & \pass & Left/right vectors have pairing one; full contractions, Taylor half-factor and opposite physical-$g$ coefficient are correct.\\
155--159 & \pass & Scalar expansion has cubic coefficient $-(A+b)$ and quadratic derivative $2b$.\\
160--169 & \pass & Original-field mixed derivatives and all five second directional contributions are correct; averaging introduces no extra $n$.\\
170--175 & \pass & Generic theorem title is justified by the ensuing full-system proof.\\
176--189 & \pass & Branch completeness, smoothness and expansions have fixed-$b$ local scope; both slope orientations are correct.\\
190--206 & \pass & Tangents are in the synchronous state--parameter plane; full nonlinear stability table and transcritical conclusion have the stated generic hypotheses.\\
207--215 & \pass & Second IFT constructs the unique small root; nonzero iff $\mu\ne0$; the other root stays excluded for fixed $b\ne0$.\\
216--225 & \pass & First/second derivatives and tangent vectors are correct; $d\mu/dg=-1$ is explicitly used.\\
226--235 & \pass & Invariant line is a fixed-parameter center manifold; its product with parameter is the full extended center manifold because $F_\mu(0,0)=0$.\\
236--244 & \pass & Coordinate $u=Ar^2-br$ is locally invertible; exact transformed vector field and rational time factor are correct.\\
245--249 & \pass & Denominators stay nonzero, time factor is positive, and orbital equivalence is limited to the center restriction.\\
250--254 & \pass & Full derivative of the scalar field is taken before imposing equilibrium; product-rule terms are correct.\\
255--261 & \pass & Uniform eigenvalue and alternative expression use the equilibrium relation with the correct minus sign.\\
262--270 & \pass & Both $b$ factors retain their sign locally; the exact sign proof is primary and the expansion only corroborates it.\\
271--282 & \pass & General P04 coefficients and phase signs remain valid at negative synchronous states and colliding offsets.\\
283--293 & \pass & Finite-label continuity yields one common strict transverse gap, including $b=0$; no selection theorem is used.\\
294--303 & \pass & Origin sign and nonlinear linearization results prove the exchange. Exactly one unstable real direction refers to the off-critical hyperbolic equilibria just treated.\\
304--315 & \pass & Physical corollary has correct branch side and stability for each sign of $b$; orthant statements are off-critical.\\
316--324 & \pass & Positive synchronous escape, with a fixed radius for each fixed $\mu>0$, supplies an admissible trajectory proving relative instability.\\
325--330 & \pass & Nonnegative invariance restricts stability; interior positive-branch instability and the separate away-from-zero root are correctly scoped.\\
331--338 & \pass & Generic critical instability is in $\U$; the $b<0$ orthant conclusion is expressly withheld.\\
339--343 & \pass & Critical synchronous equation is exact.\\
344--353 & \pass & Magnitude bound $|b|r^2/3$, sign choice for $r$, finite exit, domain preservation until exit and orthant limitation are correct.\\
354--359 & \pass & Degenerate subsection/proposition opening does not assume genericity.\\
360--367 & \pass & For $\mu<0$ and $\mu=0$, local root counts and stability types are correct; triple zero describes the scalar field.\\
368--377 & \pass & Two roots for $\mu>0$ are in $\U$ locally; both are exponentially stable there; origin is unstable and only the positive root is nonnegative.\\
378--382 & \pass & Translation to $g$ is correct; pitchfork and two-parameter unfolding are explicitly not asserted.\\
383--388 & \pass & The transverse reduction applies without $b\ne0$; both roots can be placed in a common local neighborhood with $r>-1$.\\
389--393 & \pass & Critical cubic derivative is $-6A$; scalar zero multiplicity is exactly three.\\
394--403 & \pass & Uniform eigenvalue is negative on both signed branches; transverse continuity and the earlier positive escape supply full off-critical conclusions.\\
404--415 & \pass & The fixed-parameter center-stability theorem is stated with and satisfies all hypotheses; the neutral parameter direction is explicitly excluded.\\
416--418 & \pass & Independent direct Lyapunov proof is at the fixed degenerate tuple in real mean/transverse coordinates.\\
419--426 & \pass & Circulant column sums kill all linear-in-$y$ mean terms for every $r$; uniform transverse Taylor remainder is $O(\norm y_2^2)$.\\
427--432 & \pass & Vanishing transverse field on synchrony gives the exact integral factorization and the displayed operator $B:E\to E$.\\
433--436 & \pass & Operator norm of $B$ shrinks to zero; the quadratic-form estimate follows from unitary Fourier diagonalization, or explicitly its real symmetric part.\\
437--444 & \pass & Radius choice is feasible with a fixed $K$; the $rR$ bound and denominator comparison give the strict negative derivative for both signs of $r$.\\
445--450 & \pass & Compact sublevels lie in $\U$ and the estimate neighborhood, are invariant, and ensure continuation and full asymptotic stability.\\
451--460 & \pass & Sign preservation and convergence justify the reciprocal-square derivative limit, algebraic asymptotic, and impossibility of full exponential stability.\\
461--471 & \pass & Codimension proposition concerns the positive physical parameter space and an algebraic meeting of loci only.\\
472--480 & \pass & Gradients have rank two, the regular-level theorem applies, and the locus is nonempty.\\
481--483 & \pass & The center zero remains spectrally simple; no double-zero claim.\\
484--492 & \pass & P02 formulas give equality of loci iff $a=A$, at root zero; positive repeated root and saddle-node scope are correctly distinguished.\\
493--500 & \pass & B02 resolution is limited to P09; NB07, frozen supporting role and exclusion of unfolding are respected.\\
501--512 & \pass & Analytical checks are actually in the proof; computational provenance is separate and has been reproduced, without treating validation as proof.\\
513--518 & \pass & Audit-priority preamble and IFT/slope priorities accurately identify the proof obligations.\\
519--524 & \pass & Coefficient and center-normal-form audit requests preserve normalization and full-system completeness.\\
525--528 & \pass & Spectral and physical-side priorities explicitly include negative states, collisions and open-domain versus orthant distinctions.\\
529--534 & \pass & Degenerate Lyapunov and codimension priorities are correctly scoped; no named pitchfork/cusp classification is asserted.\\
535--539 & \pass & Final exclusions retain the P10 boundary and all other forbidden later-module claims; closing environment is complete.\\
\bottomrule
\end{longtable}
% A09-LINE-LEDGER-END

\section{Requested-check coverage}
The request's checks 0--1 are addressed in Sections~\ref{sec:auth}--\ref{sec:deps};
checks 2--9 in Sections~\ref{sec:reduction} and the verdict table;
checks 10--16 in Section~\ref{sec:generic}; checks 17--25 in
Section~\ref{sec:stability}; checks 26--27 in Section~\ref{sec:critical};
checks 28--34 in Section~\ref{sec:degenerate}; checks 35--41 in
Section~\ref{sec:lyap}; checks 42--47 in Section~\ref{sec:scope}; checks
48--49 in Section~\ref{sec:checks}; and check 50 in Section~\ref{sec:lines}.
Checks 51--55 are implemented by the verdicts, exact-slug source, audit-only
archive and closure handoff. None of these dispositions depends on compiling
this document or on a developer validation result.

\section{Closure handoff to the original P09 development chat}
\label{sec:handoff}
\textbf{OVERALL VERDICT: VERIFIED. Required corrections: none.}
There are no E01, E02 or other required issue IDs to close. The original
P09 development chat may perform administrative closure without changing
the mathematical body of the exact audited P09 v1.0. Preserve the original
audit source and its evidence, and record the current/audited hashes and
verdicts in the appropriate lifecycle metadata there.

This audit chat has not edited P09, promoted it to VERIFIED in any project
file, begun P10, or created a replacement full project. The audit delivery
is named exactly \path{cyclic_network_project_A09_audit_only.zip}, with
the root \path{cyclic_network_project/} and only A09 artifacts. The
original \path{cyclic_network_project_P09_draft.zip} remains unchanged.

\textbf{After successful correction/closure, the original P09 development
chat must create the COMPLETE full-project archive:}
\begin{center}
\path{cyclic_network_project_P09_audited.zip}.
\end{center}
That archive must not be created in this audit chat.

\textbf{P10 gate: cleared after administrative A09/P09 audit closure.}
All mathematical prerequisites P02, P03, P04 and the audited P09 result
are available; P10 is the next frozen attempt. This is permission to begin
its scoped proof work after closure, not a P10 theorem verdict. P10 must
state its own positive-equilibrium, full-spectrum and local stable-coexistence
hypotheses, retain the B03/local-versus-global boundary, and undergo its own
audit. No global basin result, unfolding or later proof is authorized or
supplied by this audit.
\end{document}
